\section{Introduction}
\label{sec:introduction}

\ac{ai} systems increasingly support the production and interpretation of scientific results. Their outputs are commonly described as predictive, reproducible, explanatory, mechanistic, causal, or representative of a digital twin. These terms express different evidential commitments. A model may predict held-out observations while failing to identify the process that generated them. An analysis may execute reproducibly while relying on a non-unique explanation. A local observational association may be stable while providing no intervention-based support. The problem addressed in this work is therefore not the absence of computational results. It is the absence of an explicit and executable boundary between those results and the wording used to communicate them.

Mouse-brain modelling makes this problem particularly visible. Public atlases, tracer data, electrophysiological recordings, connectomic reconstructions, neural-response benchmarks, and simulation platforms describe different biological scales \citep{oh2014mesoscale,melozzi2017virtualmouse,dai2020bmtk,billeh2020v1,willeke2022sensorium,bae2025microns}. Static or dynamic neural-response prediction is not evidence that a synaptic topology has been identified. A measured synapse and a functional descriptor in a local cortical volume do not establish a whole-brain causal model. A brain-wide simulation based on inferred connectivity does not become an entity-specific digital twin solely because a macroscopic signal is reproduced. These distinctions remain relevant even when every dataset and program is publicly available.

Several mature areas address parts of the problem. \ac{vv} and \ac{uq} connect model credibility to a stated context of use \citep{oberkampf2010verification,riedmaier2021vvuq}. Assurance cases decompose broad claims into arguments and evidence \citep{mohamad2021assurance,porter2024praise}. Scientific claim verification connects natural-language statements to documentary evidence \citep{wadden2020scifact,vladika2023survey}. Reproducibility programmes improve the traceability of code, data, and execution \citep{pineau2021reproducibility,barker2022fair4rs}. These contributions are necessary, but they do not by themselves determine which interpretation of a heterogeneous computational result is admissible.

The resulting task is a decision-support problem. For each candidate statement, an author or reviewer must decide whether to retain it, weaken it, block it, defer it, or request judgement that cannot be automated. The criteria are non-interchangeable. Strong predictive performance cannot compensate for an absent intervention. Repeated execution cannot compensate for a failed topology control. This structure is related to non-compensatory multicriteria decision making, where a veto prevents strong criteria from hiding a decisive weakness \citep{figueira2013electre}. The object of the decision is different here. It is the admissible scope of one scientific statement.

This paper presents MouseClaimBench, an executable \ac{dss} for computational mouse-brain claims. A domain contract specifies the evidence blocks required by a claim, the rule applied to each block, the source artifact, and the wording that remains admissible. The system retains every observation in its original scale. It does not map a correlation, a bootstrap interval, and a reliability coefficient to a common synthetic score. It also distinguishes computational reproducibility, internal reproduction within one resource, and external replication. These distinctions remove two ambiguities in the earlier design and make the decision artifact more suitable for scientific review.

The evaluation is organized to avoid treating software self-consistency as independent validity. First, a new oracle benchmark generates five \ac{sem} regimes whose graph and intervention status are fixed before the evidence rules are executed. It produces finite-sample errors rather than a perfect result defined by construction. Second, four real mouse-brain cases are rebuilt from observed fields in Allen \ac{vbn}, Static Sensorium, Dynamic Sensorium, and \ac{microns} artifacts. Third, SciFact and the Tuebingen Cause-Effect Pairs expose retrieval and directional-overclaiming failures outside the handcrafted contract cases \citep{wadden2020scifact,mooij2016cause}. The original 144-case benchmark is retained only as a software-conformance test because its labels share the semantics of the legacy gate.

The results show a clear policy trade-off. The non-compensatory contract reduces the \ac{fpr} from 0.271 for a prediction-only shortcut to 0.019, while its \ac{fnr} increases from 0.007 to 0.070. In the paired comparison, the contract alone is correct for 747 decisions and the shortcut alone is correct for 149 decisions. The real-data matrix authorizes bounded prediction, internal reproduction, topology, and local structure--function claims. It does not authorize causality, external biological replication, or a complete mouse-brain digital twin. These are useful negative decisions because they identify the exact evidence that remains absent.

The contribution is methodological and implementational. MouseClaimBench does not introduce a new neural predictor, simulator, causal-discovery algorithm, or textual fact checker. It introduces an inspectable decision policy, applies it to real mouse-brain evidence, and evaluates its behaviour against independently generated structural truth and external failure modes. This scope aligns with the foundations, functionality, implementation, and evaluation concerns of \acp{dss}, while leaving impact on human decision makers as an explicit untested question \citep{arnott2005critical,hevner2004design,peffers2007design}.

The remainder of the paper is organized as follows. Section~\ref{sec:contributions} states the contributions and boundaries. Section~\ref{sec:related-work} positions the framework against decision support, assurance, claim verification, and mouse-brain modelling. Section~\ref{sec:method} defines the contract and evidence states. Section~\ref{sec:experiments} presents the protocol and results. Section~\ref{sec:discussion} examines interpretation, threats to validity, and the pending human study. Section~\ref{sec:conclusions} summarizes the findings.
